\label{sec:taylor}

\begin{exposition}
Taylor's theorem is the derivative's most ambitious claim: local behaviour to
finite order is encoded by derivatives at one point, with an error term that
can be estimated. The first-order case is the seed of the closing reframe,
where the derivative stops being only a number \(f'(c)\) and becomes the
linear map \(h\mapsto f'(c)h\).
\end{exposition}

% ---------------------------------------------------------
% Definition --- Taylor Polynomial at a Point
% ---------------------------------------------------------

\begin{definitionbox}{Definition (Taylor Polynomial at a Point)}
\begin{definition}[Taylor Polynomial at a Point]
\label{def:taylor-polynomial-at-a-point}
Let $I\subseteq\mathbb{R}$ be an interval, let $a\in I$, let
$n\in\mathbb{N}$, and let $f:I\to\mathbb{R}$ have derivatives
$f^{(k)}(a)$ for $0\le k\le n$. The \textbf{Taylor polynomial of order}
$n$ for $f$ at $a$ is
\[
  T_{n,a}f(x)
  :=
  \sum_{k=0}^{n}\frac{f^{(k)}(a)}{k!}(x-a)^k.
\]
\end{definition}
\end{definitionbox}

\begin{remark*}[Standard quantified statement]
\[
\begin{aligned}
&\text{Let } I\subseteq\mathbb{R},\; a\in I,\; n\in\mathbb{N},\;
f:I\to\mathbb{R}.\\
&\text{If } f^{(k)}(a)\text{ exists for every }0\le k\le n,\text{ then}\\
&\qquad
\forall x\in I\;
\left(
  T_{n,a}f(x)=\sum_{k=0}^{n}\frac{f^{(k)}(a)}{k!}(x-a)^k
\right).
\end{aligned}
\]
\end{remark*}

\begin{remark*}[Interpretation]
The Taylor polynomial is the unique polynomial of degree at most $n$ whose
first $n$ derivatives at $a$ agree with those of $f$. Its coefficients are
not chosen by fitting sample points; they are dictated by derivative data at
one point.
\end{remark*}

\begin{dependencies}
\begin{itemize}
  \item \hyperref[def:derivative-at-a-point]{Derivative at a Point}
  \item \hyperref[cor:power-rule-special-case]{Power Rule for Integer Powers of a Function}
  \item \hyperref[cor:finite-sum-rule]{Finite Sum Rule}
\end{itemize}
\end{dependencies}

% ---------------------------------------------------------
% Definition --- Taylor Remainder
% ---------------------------------------------------------

\begin{definitionbox}{Definition (Taylor Remainder)}
\begin{definition}[Taylor Remainder]
\label{def:taylor-remainder}
Let $I\subseteq\mathbb{R}$ be an interval, let $a\in I$, let
$n\in\mathbb{N}$, and let $f:I\to\mathbb{R}$ have a Taylor polynomial
$T_{n,a}f$. The \textbf{Taylor remainder of order} $n$ at $a$ is the
function $R_{n,a}f:I\to\mathbb{R}$ defined by
\[
  R_{n,a}f(x):=f(x)-T_{n,a}f(x).
\]
\end{definition}
\end{definitionbox}

\begin{remark*}[Standard quantified statement]
\[
\forall x\in I\;\bigl(R_{n,a}f(x)=f(x)-T_{n,a}f(x)\bigr).
\]
\end{remark*}

\begin{remark*}[Interpretation]
The remainder is the exact error made by replacing $f$ with its Taylor
polynomial. Taylor's theorem is valuable because it gives usable formulas
or estimates for this error.
\end{remark*}

\begin{dependencies}
\begin{itemize}
  \item \hyperref[def:taylor-polynomial-at-a-point]{Taylor Polynomial at a Point}
\end{itemize}
\end{dependencies}

% ---------------------------------------------------------
% Definition --- Maclaurin Polynomial
% ---------------------------------------------------------

\begin{definitionbox}{Definition (Maclaurin Polynomial)}
\begin{definition}[Maclaurin Polynomial]
\label{def:maclaurin-polynomial}
Let $n\in\mathbb{N}$, and let $f$ be defined on an interval containing
$0$ with derivatives $f^{(k)}(0)$ for $0\le k\le n$. The
\textbf{Maclaurin polynomial of order} $n$ for $f$ is
\[
  T_{n,0}f(x)
  :=
  \sum_{k=0}^{n}\frac{f^{(k)}(0)}{k!}x^k.
\]
\end{definition}
\end{definitionbox}

\begin{remark*}[Standard quantified statement]
\[
\forall x\in I\;
\left(
  T_{n,0}f(x)=\sum_{k=0}^{n}\frac{f^{(k)}(0)}{k!}x^k
\right).
\]
\end{remark*}

\begin{remark*}[Interpretation]
Maclaurin polynomials are Taylor polynomials centered at the origin. They are
not a separate approximation theory; they are the centered-at-zero case of
the same construction.
\end{remark*}

\begin{dependencies}
\begin{itemize}
  \item \hyperref[def:taylor-polynomial-at-a-point]{Taylor Polynomial at a Point}
\end{itemize}
\end{dependencies}

% ---------------------------------------------------------
% Theorem --- Taylor's Theorem with Lagrange Remainder
% ---------------------------------------------------------

\begin{theorembox}{Theorem (Taylor's Theorem with Lagrange Remainder)}
\begin{theorem}[Taylor's Theorem with Lagrange Remainder]
\label{thm:taylor-theorem-lagrange-remainder}
Let $a,b\in\mathbb{R}$ with $a<b$, let $n\in\mathbb{N}$, and let
$f:[a,b]\to\mathbb{R}$. Suppose that $f,f',\ldots,f^{(n)}$ are continuous
on $[a,b]$ and that $f^{(n+1)}$ exists on $(a,b)$. Then, for every
$x\in(a,b)$, there exists $c$ strictly between $a$ and $x$ such that
\[
  f(x)
  =
  \sum_{k=0}^{n}\frac{f^{(k)}(a)}{k!}(x-a)^k
  +
  \frac{f^{(n+1)}(c)}{(n+1)!}(x-a)^{n+1}.
\]

\smallskip
\noindent
\hyperref[prf:taylor-theorem-lagrange-remainder]{\textit{Go to proof.}}
\end{theorem}
\end{theorembox}

\begin{remark*}[Standard quantified statement]
\[
\begin{aligned}
&\forall x\in(a,b)\;\exists c\in(a,x)\\
&\quad
\left(
  f(x)
  =
  \sum_{k=0}^{n}\frac{f^{(k)}(a)}{k!}(x-a)^k
  +
  \frac{f^{(n+1)}(c)}{(n+1)!}(x-a)^{n+1}
\right).
\end{aligned}
\]
\end{remark*}

\begin{remark*}[Predicate reading]
The displayed quantified statement is the predicate-level form of $label: its hypotheses name the input conditions, and its conclusion names the asserted structure or relation.
\[
  f\colon A\to\mathbb{R},\\
  \operatorname{Derivative}(D,f,a,\mathbb{R},\mathbb{R}).
\]
\end{remark*}

\begin{remark*}[Interpretation]
Taylor's theorem says that the error after truncating at degree $n$ has the
same shape as the next Taylor term, but with the next derivative evaluated at
some intermediate point rather than at the center. The point $c$ is generally
not known explicitly; the theorem is most often used to estimate the size of
the remainder.
\end{remark*}

\begin{dependencies}
\begin{itemize}
  \item \hyperref[def:taylor-polynomial-at-a-point]{Taylor Polynomial at a Point}
  \item \hyperref[def:taylor-remainder]{Taylor Remainder}
  \item \hyperref[thm:mean-value-theorem]{Mean Value Theorem}
  \item \hyperref[cor:finite-sum-rule]{Finite Sum Rule}
  \item \hyperref[cor:power-rule-special-case]{Power Rule for Integer Powers of a Function}
\end{itemize}
\end{dependencies}

% ---------------------------------------------------------
% Corollary --- Taylor Expansion with Peano Remainder
% ---------------------------------------------------------

\begin{corollarybox}{Corollary (Taylor Expansion with Peano Remainder)}
\begin{corollary}[Taylor Expansion with Peano Remainder]
\label{cor:taylor-expansion-peano-remainder}
Let $I\subseteq\mathbb{R}$ be an interval, let $a\in I$, let
$n\in\mathbb{N}$, and let $f:I\to\mathbb{R}$ have derivatives up to order
$n$ in a neighbourhood of $a$. If $f^{(n)}$ is continuous at $a$, then
\[
  f(x)
  =
  \sum_{k=0}^{n}\frac{f^{(k)}(a)}{k!}(x-a)^k
  +
  o\bigl((x-a)^n\bigr)
  \qquad\text{as }x\to a.
\]

\smallskip
\noindent
\hyperref[prf:taylor-expansion-peano-remainder]{\textit{Go to proof.}}
\end{corollary}
\end{corollarybox}

\begin{remark*}[Standard quantified statement]
\[
\begin{aligned}
&\lim_{x\to a}
\frac{
  f(x)-\sum_{k=0}^{n}\frac{f^{(k)}(a)}{k!}(x-a)^k
}{
  (x-a)^n
}
=0.
\end{aligned}
\]
\end{remark*}

\begin{remark*}[Interpretation]
The Peano form is a local asymptotic statement. It says that the Taylor
polynomial captures all terms through order $n$, and the remaining error is
smaller than $(x-a)^n$ near the center.
\end{remark*}

\begin{dependencies}
\begin{itemize}
  \item \hyperref[def:taylor-polynomial-at-a-point]{Taylor Polynomial at a Point}
  \item \hyperref[def:taylor-remainder]{Taylor Remainder}
  \item \hyperref[def:limit-function]{Limit of a Function}
  \item \hyperref[def:derivative-at-a-point]{Derivative at a Point}
\end{itemize}
\end{dependencies}

% ---------------------------------------------------------

\begin{corollarybox}{Corollary (First-Order Peano Remainder)}
\begin{corollary}[First-Order Peano Remainder]
\label{cor:first-order-peano-remainder}
Let $f$ be differentiable at $c$. Then, as $h\to 0$,
\[
  f(c+h)=f(c)+f'(c)h+o(h).
\]
\hyperref[prf:first-order-peano-remainder]{Go to proof.}
\end{corollary}
\end{corollarybox}

\begin{remark*}[Standard quantified statement]
\[
\operatorname{IsDifferentiable}(f,c,\mathbb{R},\mathbb{R})
\;\Rightarrow\;
\lim_{h\to 0}
\frac{f(c+h)-f(c)-f'(c)h}{h}
=0.
\]
\end{remark*}

\begin{remark*}[Interpretation]
The first-order Peano formula is the derivative rewritten as an approximation theorem. It says that the affine approximation $f(c)+f'(c)h$ captures all first-order change, and the remaining error is smaller than $h$ itself.
\end{remark*}

\begin{dependencies}
\begin{itemize}
  \item \hyperref[def:derivative-at-a-point]{Derivative at a Point}
\end{itemize}
\end{dependencies}

\begin{exposition}
\phantomsection\label{ex:sinx-limit}
Taylor expansion gives
\[
  \sin x=x-\frac{x^3}{6}+o(x^3),
\]
so
\[
  \frac{\sin x-x}{x^3}
  =
  \frac{-x^3/6+o(x^3)}{x^3}
  \longrightarrow
  -\frac16.
\]
L'H\^{o}pital's rule reaches the same value by three differentiations:
\[
  \lim_{x\to0}\frac{\sin x-x}{x^3}
  =
  \lim_{x\to0}\frac{\cos x-1}{3x^2}
  =
  \lim_{x\to0}\frac{-\sin x}{6x}
  =
  \lim_{x\to0}\frac{-\cos x}{6}
  =
  -\frac16.
\]
Taylor exposes the reason for the answer: the limit is the cubic coefficient.
\end{exposition}

% ---------------------------------------------------------

\begin{propositionbox}{Proposition (The Flat Function: Smooth but Not Analytic)}
\begin{proposition}[The Flat Function: Smooth but Not Analytic]
\label{prop:flat-function}
Define \(f:\mathbb{R}\to\mathbb{R}\) by
\[
  f(x)=
  \begin{cases}
    e^{-1/x^2}, & x\neq 0,\\
    0, & x=0.
  \end{cases}
\]
Then:
\begin{enumerate}
  \item \(f\in C^\infty(\mathbb{R})\);
  \item \(f^{(n)}(0)=0\) for every \(n\geq 0\);
  \item the Maclaurin series of \(f\) is identically \(0\), so \(f\notin C^\omega\) on any neighbourhood of \(0\).
\end{enumerate}
In particular, \(C^\infty(\mathbb{R})\setminus C^\omega(\mathbb{R})\neq\varnothing\).
\hyperref[prf:flat-function]{Go to proof.}
\end{proposition}
\end{propositionbox}

\begin{remark*}[Standard quantified statement]
\[
  \left(
  \forall n\in\mathbb{N}: f^{(n)}(0)=0
  \right)
  \land
  \left(
  \forall r>0\;\exists x:\ 0<|x|<r\land f(x)\neq 0
  \right)
  \Longrightarrow
  f\in C^\infty\setminus C^\omega.
\]
\end{remark*}

\begin{remark*}[Predicate reading]
The displayed quantified statement is the predicate-level form of $label: its hypotheses name the input conditions, and its conclusion names the asserted structure or relation.
\[
  f\colon A\to\mathbb{R},\\
  \operatorname{Derivative}(D,f,a,\mathbb{R},\mathbb{R}).
\]
\end{remark*}

\begin{remark*}[Interpretation]
The Taylor series sees only derivative data at the center, and all of that data is zero for the flat function at \(0\). The function is nevertheless positive away from \(0\). This separates smoothness from analyticity and supplies the strictness witness referenced by the smoothness tower.
\end{remark*}

\begin{dependencies}
\begin{itemize}
  \item \hyperref[def:higher-derivatives]{Definition: Higher Derivatives}
  \item \hyperref[def:class-cinfty]{Definition: Smooth Function; Class \(C^\infty\)}
  \item \hyperref[def:class-comega]{Definition: Real-Analytic Function; Class \(C^\omega\)}
  \item \hyperref[def:maclaurin-polynomial]{Definition: Maclaurin Polynomial}
  \item \hyperref[thm:smoothness-tower]{Theorem: The Smoothness Tower}
\end{itemize}
\end{dependencies}
